\section{Ricerca nello spazio degli stati}

%

\section{Valutazione delle strategie}

%

\section{Strategie di ricerca non informate}

%

\section{Ricerca in ambienti complessi}

%

\section{Soddisfacimento dei vincoli}








\begin{comment}









Durante la risoluzione del problema si calcola il path cost, la sommatoria dei path step richiesti per arrivare alla soluzione. Naturalmente si dovrà scegliere il costo minore. Notare inoltre che è comune applicare una noOP, ovvero il non fare operazioni né cambiare stato. È importante implementarla perché alle volte spostarsi di stato è più costoso rispetto allo stare fermi, dato il path step cost.

--- Tree-search problem
Fare un'esplorazione simulata e offline. Si pianifica una sequenza, insomma. Si generano i successori degli stati già osservati, ovvero si espandono gli stati.

```
function Tree-Search(problem, strategy) returns solution or failure
    initialize search tree using initial state of problem

    do
        if no candidates for expansion then return failure
        choose a leaf node for expansion according to strategy
        if node contains a goal state then return solution
        else add successor nodes to search tree
    end
```

L'algoritmo di ricerca degli alberi è BFS/DFS. Uno stato è una configurazione fisica del nodo, il quale contiene strutture dati con vari dati, come path cost e simili.

In merito all'esplorazione, viene creata la frontiera: struttura dati dove s pongono i nodi in un certo ordine, determinato dalla strategia applicata.
La frontiera ha quindi la sequenza dei passi da fare e se termina con una soluzione, si avrà avuto successo, altrimenti riporta un fallimento.

All'interno della frontiera si effettua un pop dalla frontiera. Praticamente è una coda LIFO. Lo pseudocodice generale è:

```
function Tree-Search (problem, frontier) return solution or failure
    frontier \gets insert(Make-Node(problem.initial-State))

    while not isEmpty(frontier) do
        node \gets Pop(frontier)
        if problem.Goal-Test(node.State) then return node
        frontier \gets InsertAll(Expand(node, problem))
    end loop
    return failure
```

Poi c'è il codice relativo all'espansione, che ritorna un insieme di nodi successori:
```
function Expand(node, problem) returns a set of nodes
    successors \gets emptyset

    for each action, result in Successor-Fn(problem, node.State)
        do
            s \gets new Node
            s.Parent-Node \gets Node
            s.Action \gets action
            s.State \gets result
            s.Path-Cost \gets node.Path-Cost + Step-Cost(node.State, action, result)
            s.Depth \gets node.Depth+1
            add s to successors set
        return successors
```

Come già detto, le strategie sono valutate in più ambiti:
    - completezza: Trova sempre una soluzione se esiste
    - complessità temporale: Numero di nodi generati/espansi
    - Complessità spaziale: Massimo numero di nodi in memoria
    - Ottimalità: Trova sempre la soluzione di costo minimo

Le complessità temporali e spaziali sono misurate in termini di:
    - b: Massimo numero di figli che un nodo può avere all'interno dell'albero di ricerca
    - d: profondità della soluzione di costo minimo
    - m: profondità massima dello spazio degli stati (possibile che sia infinita)


--- Problema degli stati ripetuti
Non scoprire stati ripetuti potrebbe far diventare un problema di complessità lineare in uno esponenziale. Non ideale. Ciò è evitabile aggiungendo un insieme di stati già esplorati, simile a come si fa per gli alberi rb.

La funzione in pseudocodice che implementa questa caratteristica è la graph search:
```
function Graph-Search(problem, frontier) return solution or failure
    explored \gets emptyset
    frontier \gets insert(Make-Node(problem.Initial-State))

    while not IsEmpty(frontier) do
        node \gets Pop(frontier)
        if problem.Goal-Test(node.State) then return node
        if node.State is not in explored then
            add node.State to explored
            frontier \gets InsertAll(Expand(node, problem))
        end if
    end loop
    return failure
```




    Lez3 Intelligenza artificiale 08-10-26

--- General tree search
Agenti che risolvono problemi tramite la ricerca

Tree search algorithm consente di trovare una soluzione negli stati interni all'agente.
Differenza tra stato e nodo

Abbiamo un problema (conoscenza dello stato iniziale e dell'environment) e una frontiera che contiene dati.

```
function Tree-search (problem, frontier) returns solution or failure
frontier <- insert(make-node(problem.initialState))
while not isEmpty(frontier) do
node <- pop(frontier)

... Manca il resto
```

--- Expand nodes
function expand(node, problem) returns a set of nodes

```
successors <- empty set
for each action, result in Successor-Fn(problem, node.state) do
	s<- new node
	s<- Parent-node node
	s.Action <- 	
... manca il resto
```

--- Strategie di ricerca
Una strategia è definita prendendo l'ordine di espansione dei nodi e si valutano nelle dimensioni di correttezza e completezza (garantire sempre una soluzione se esiste)
complessità di tempo e spazio, come anche di ottimalità (riesce a riportare sempre la soluzione migliore?). Complessità su tempo e spazio si misurano con:
	- b: maximum branching factor: Numero massimo di figli che ciascun nodo ha
	- d: profondità della soluzione a costo minimo
	- m: Profondità massima dell'albero, quindi dello state space. Potrebbe essere infinito.
	
- stati ripetuti
Non notare stati ripetuti può portare ad un espansione di un problema lineare in uno esponenziale: node^d, in pratica.

--- Graph search
Introduce una caratteristica ai nodi: explored, che indica se un nodo è già stato esplorato o meno. Consente di saltare quindi nodi già esplorati.

```
inserisci algoritmo graph search
```



----- Esempio ricerca problema vacuum cleaner
State: <Locazione, r_a State, r_b State> \equiv <(A, B), (C, D), (C, D)>
Actions: <R, L, S>, R: A->B, L: B->A, S: D->C, S: C->C

Il goal è che tutte le piastrelle siano C, ovvero clean.
Il problema è fully observable, lo sviluppo è, con tree-search:

<A, C, D>

Se faccio S genero <A, C, D>
	S genera ora <A, C, D>
	R genera ora <B, C, D>

Se faccio R genero <B, C, D>
	L genera <A, C, D>
	S genera <B, C, C>	// Questo è un nodo di goal
	
b = 2	
d = 2
m = \infty
	
Tuttavia i nodi di goal devono essere estratti dalla frontiera espansa.
Generalmente non si espande mai uno stesso nodo nello stesso branch, perché normalmente non cambia lo stato.
\end{comment}

\begin{comment}
    --- Strategie di ricerca non informate
Approcci di ricerca senza sapere nulla in merito al problema, perché usano solo le informazioni proposte nella definizione del problema. Sono attuabili con alcuni algoritmi come
	- BFS graph search: Frontiera fifo ed espande il nodo più superficiale non espanso. TODO Aggiungi algoritmo
		Proprietà: è completa se b è finito, compl tempo: O(b^d), compl. spazio: O(b^d), è ottima se si ha il medesimo costo per ogni step nell'albero, non in generale.
	- Uniform-cost search: Espande il nodo non espanso dal costo minimo. Ordiniamo la frontiera per posto tenendo in mente sempre il nodo che costa meno.
		Se i costi di spostamento sono uguali, ha proprietà equivalenti a BFS graph search.
		Due modifiche da apportare per garantire ottimalità sempre:
			1. Esegui il goal test when selecting nodes for expansion and not when generated.
			2. Check if child node is already present inside the frontier with a higher cost, then replace the previous node.
	- DFS graph search: Frontiera è una coda Lifo.
		Non è completa, fallisce in loop infiniti. Tuttavia si può modificare per eludere
		stati ripetuti. E' tuttavia completa in spazi finiti.
		Tempo: O(b^m), atroce se m è molto più grande rispetto a d, ma se le soluzioni sono dense
		potrebbe essere più veloce di BFS. Spazio: O(bm). Non è ottimale, né è possibile renderla tale.
	 - Iterative deepening search: Esegui una DFS definendo un limite alla profondità dell'albero.
		Il senso di questa strategia, è limitare la complessità spaziale. Introduce poi la possibilità di ritornare successo, fallimento e cutoff, in base al threshold scelto.
		Si tratta di una soluzione completa, tempo O(b^d), spazio O(bd). Inoltre, se gli step costano 1 può essere modificata per esplorare alberi di costo uniforme.

--- Wolf Sheep Cabbage problem
Environment: <W, S, C, M, B>
Init: <0,0,0,0,0>, Goal: <1,1,1,1,1>

Action		Pre-cond				From					To
carryW		s\neq c, w=m, m=b		(w,s,c,m,d) 	-> 		(!w,s,c,!m,!b)
carryS		m=b, s=m				(w,s,c,m,d)		->		(w,!s,c,!m,!b)
carryC		m=b, m=c, w\neq s		(w,s,c,m,d)		->		(w,s,!c,!m,!b)
carryNthg	m=b, w\neq s, n\neq c	(w,s,c,m,d)		->		(w,s,c,!m,!b)

Notare che m=b sempre, in ogni azione. E' quindi possibile semplificare rimuovendo
questo elemento dalla tupla in toto, semplificando lo stato.
\end{comment}